% Spectral reflection symmetry and periodic multiplicity structure
% for quantum rotors on SO(3)/H
% Target: J. Phys. A: Mathematical and Theoretical
% v10 — Feb 2026

\documentclass[12pt]{article}

% ============================================================
% Packages
% ============================================================
\usepackage{amsmath, amssymb, amsthm}
\usepackage{mathtools}
\usepackage{booktabs}
\usepackage{hyperref}
\usepackage[utf8]{inputenc}

% ============================================================
% Theorem environments
% ============================================================
\newtheorem{theorem}{Theorem}[section]
\newtheorem{corollary}[theorem]{Corollary}
\newtheorem{lemma}[theorem]{Lemma}
\newtheorem{proposition}[theorem]{Proposition}
\theoremstyle{definition}
\newtheorem{definition}[theorem]{Definition}
\theoremstyle{remark}
\newtheorem{remark}[theorem]{Remark}

% ============================================================
% Notation
% ============================================================
\newcommand{\ZZ}{\mathbb{Z}}
\newcommand{\CC}{\mathbb{C}}
\newcommand{\RR}{\mathbb{R}}
\newcommand{\Hil}{\mathcal{H}}
\newcommand{\identity}{\mathbf{1}}
\newcommand{\SU}{\mathrm{SU}}
\newcommand{\SO}{\mathrm{SO}}
\newcommand{\UU}{\mathrm{U}}
\newcommand{\Aut}{\mathrm{Aut}}
\DeclareMathOperator{\Hom}{Hom}
\DeclareMathOperator{\ord}{ord}
\DeclareMathOperator{\lcm}{lcm}

% ============================================================
% Title
% ============================================================
\title{Spectral reflection symmetry and periodic multiplicity structure\\
for quantum rotors on $\SO(3)/H$}
\author{}  % TODO: add authors
\date{}

\begin{document}
\maketitle

% ============================================================
\begin{abstract}
The scalar spectrum of a quantum rotor on $\SO(3)/H$
decomposes into superselection sectors labelled by
one-dimensional characters of $\pi_1 = 2H \subset \SU(2)$.
We prove that the sub-threshold multiplicities satisfy the
reflection identity
$m(\chi,j) + m(\chi, j^*{-}j) = 1$
for every scalar character~$\chi$
if and only if a divisibility condition~(B) on element orders
holds, and that~(B) is equivalent to the commutator
obstruction $-1 \in [2H,2H]$.
Under~(B), the multiplicity function decomposes as a linear
drift plus a periodic binary residual with period $P = |H|/2$,
from which binary occupation, regular-representation tiling,
and threshold sharpness follow as immediate corollaries.
Under symmetry breaking $H \to K$ with $Q_8 \not\subset 2K$,
spinorial sectors reappear; for odd dihedral groups,
the reflection fails globally but persists in
spinorial sectors.
\end{abstract}

% ============================================================
\section{Introduction}
\label{sec:intro}
% ============================================================

A rigid quantum rotor whose orientational symmetry is that of a
finite subgroup $H \subset \SO(3)$~\cite{BunkerJensen2006,Zhilinskii2001}
lives on the configuration space $M = \SO(3)/H$.
Because $M$ is multiply connected, with fundamental group
$\pi_1(M) \cong 2H \subset \SU(2)$, quantization splits into
superselection sectors indexed by one-dimensional characters
$\chi \colon 2H \to \UU(1)$~\cite{Schulman1968, LaidlawDeWitt1971,
Dowker1972, Isham1984}.
The Hamiltonian $\hat{H} = \mathbf{L}^2/2I$ gives eigenvalues
$E_j = Bj(j{+}1)$, where $B = \hbar^2/2I$, and each sector
carries its own multiplicity function
$m(\chi,j)$ counting how often level~$j$ appears.
Under a $2\pi$ rotation the wave function acquires the phase
$\chi(-1)$: sectors with $\chi(-1) = +1$ are
\emph{tensorial}, those with $\chi(-1) = -1$ \emph{spinorial}.

For the polyhedral groups ($T$, $O$, $I$), the element
$-1$ lies in the commutator subgroup of~$2H$, forcing
$\chi(-1) = +1$ for every one-dimensional character:
no spinorial sector exists.
Under a trigonal distortion $O \to D_3$, the obstruction lifts:
spinorial sectors reappear with gap $3.75B$, while the
trivial-sector gap reduces from $20B$ to $6B$
(Table~\ref{tab:densification}).

\begin{theorem}[Spectral reflection $\Leftrightarrow$ divisibility]
\label{thm:main}
Let $H \subset \SO(3)$ be a finite subgroup with
$|H|$ even, and set $j^* = |H|/2 - 1$.
The following are equivalent:
\begin{enumerate}
  \item[\textup{(B)}] $|H|/\ord(h)$ is even for every
    $h \in H \setminus\{e\}$;
  \item[\textup{(C)}] $m(\chi,j) + m(\chi, j^*{-}j) = 1$
    for every scalar character~$\chi$ and every
    $j \le j^*$ of matching parity.
\end{enumerate}
The proof is given in \S\ref{sec:duality-theorem}.
\end{theorem}

\begin{proposition}[Topological link]
\label{prop:intro-A-iff-B}
For $G = 2H \subset \SU(2)$, condition~\textup{(B)} is
equivalent to~\textup{(A)}: $-1 \in [G,G]$
\textup{(}all scalar sectors are tensorial\textup{)}.
The proof is given in \S\ref{sec:topological}.
\end{proposition}

\noindent
Together, Theorem~\ref{thm:main} and
Proposition~\ref{prop:intro-A-iff-B} give the triple equivalence
\textup{(A)}$\,\Leftrightarrow\,$\textup{(B)}$\,\Leftrightarrow\,$\textup{(C)}.

Individual octahedral multiplicities are
classical~\cite{Hecht1960}, branching rules
$\SU(2) \to 2H$ are standard~\cite{BrockerTomDieck1985,AltmannHerzig1994},
and the group-theoretic equivalence
(A)$\,\Leftrightarrow\,$(B) follows from known
structure theory of finite subgroups of $\SU(2)$.
What has not been observed previously is:
\begin{enumerate}
\item[(i)] the equivalence
  (B)$\,\Leftrightarrow\,$(C): the sub-threshold
  reflection $m(\chi,j) + m(\chi,j^*{-}j) = 1$
  holds for \emph{all} scalar characters~$\chi$
  if and only if the exponent saturates at $|H|/2$;
\item[(ii)] sector-selective duality: when~(B) fails
  for odd dihedral groups, the reflection survives in
  spinorial sectors because the violating half-turn
  classes are spectrally invisible at half-integer~$j$;
\item[(iii)] the periodic staircase:
  under~(B), $m(\chi, j{+}P) = m(\chi, j) + 1$
  for $P = |H|/2$, yielding binary occupation,
  regular-representation tiling, and threshold sharpness
  as immediate corollaries.
\end{enumerate}
For the trivial character~$A_0$, the reflection
follows from the Gorenstein property of the invariant
ring $\CC[x,y]^{2H}$~\cite{Watanabe1974}; however,
Gorenstein duality does not extend to non-trivial
characters, and the full equivalence
(B)$\,\Leftrightarrow\,$(C) for all scalar~$\chi$
requires the mode-antisymmetry argument of
Theorem~\ref{thm:duality}.

Section~\ref{sec:framework} sets up the quantization framework.
Section~\ref{sec:duality} proves the spectral reflection theorem
and derives the periodic staircase.
Section~\ref{sec:obstruction} connects the divisibility condition
to the topological obstruction and examines re-entrance under
symmetry breaking.
Section~\ref{sec:discussion} discusses scope and open problems.

% ============================================================
\section{Framework}
\label{sec:framework}
% ============================================================

\subsection{Configuration space and quantization sectors}
\label{sec:config}

For a finite subgroup $H \subset \SO(3)$, the configuration
space $M = \SO(3)/H$ has fundamental group
$\pi_1(M) \cong 2H$, the binary lift of~$H$
in~$\SU(2)$~\cite{Mermin1979}, with $|2H| = 2|H|$.
Path-integral quantization on~$M$ produces inequivalent
scalar sectors indexed by one-dimensional characters
$\chi \colon 2H \to \UU(1)$~\cite{Schulman1968, LaidlawDeWitt1971,
ImboSudarshan1988, FinkelsteinRubinstein1968,
Giulini1995, BalachandranEtAl1983, Mouchet2021}.
Since one-dimensional characters factor through the
abelianization, the number of scalar sectors is
$\kappa = |(2H)^{\mathrm{ab}}|$.
Explicitly: $(2C_n)^{\mathrm{ab}} \cong \ZZ_{2n}$,
$(2D_n)^{\mathrm{ab}} \cong \ZZ_4$ ($n$ odd) or
$\ZZ_2 \times \ZZ_2$ ($n$ even),
$(2T)^{\mathrm{ab}} \cong \ZZ_3$,
$(2O)^{\mathrm{ab}} \cong \ZZ_2$,
$(2I)^{\mathrm{ab}} = \{1\}$.
We write $A_0$ for the trivial character
($A_0(g) = 1$ for all~$g$).

The rigid-rotor Hamiltonian $\hat{H} = \mathbf{L}^2/2I$ has
eigenvalues
\begin{equation}
  \label{eq:energy}
  E_j = Bj(j{+}1), \qquad B = \hbar^2/2I,
\end{equation}
with degeneracy $2j{+}1$ (magnetic substates
$m = -j, \ldots, j$).
The element $-1 \in \SU(2)$ belongs to every binary
group~$2H$ and is central; since $(-1)^2 = 1$, every
one-dimensional character satisfies $\chi(-1) = \pm 1$.
A sector is \emph{tensorial} if $\chi(-1) = +1$ and
\emph{spinorial} if $\chi(-1) = -1$.

\subsection{Mode functions and multiplicity formula}
\label{sec:modes}

The spin-$j$ representation $V_j$ of~$\SU(2)$ has
dimension $2j{+}1$ and character
\begin{equation}
  \label{eq:su2-character}
  \chi_j(g)
  = \frac{\sin\bigl((2j{+}1)\alpha\bigr)}{\sin\alpha}\,,
\end{equation}
where $\alpha$ is the half-angle of~$g$
(i.e., $g$ is conjugate to
$\mathrm{diag}(e^{i\alpha}, e^{-i\alpha})$).
The multiplicity of a scalar character $\chi$ in the
restriction $V_j|_{2H}$ is~\cite{BrockerTomDieck1985}
\begin{equation}
  \label{eq:multiplicity}
  m(\chi, j)
  = \frac{1}{|2H|}\sum_{g \in 2H}
    \overline{\chi(g)}\,\chi_j(g).
\end{equation}

Each non-central conjugacy class of~$H$ has elements of
common $\SO(3)$~order~$d$; the corresponding elements
in~$2H$ have half-angle $\alpha = \pi/d$.
Define the \emph{mode function}
\begin{equation}
  \label{eq:mode}
  f(\alpha, j)
  = \frac{\sin\bigl((2j{+}1)\alpha\bigr)}{\sin\alpha}\,.
\end{equation}
Separating the two central elements $g = \pm 1$
from~\eqref{eq:multiplicity}:
\begin{equation}
  \label{eq:mode-decomposition}
  m(\chi, j)
  = \frac{1}{|2H|}
    \Bigl[
      (2j{+}1) + \chi(-1)\,(-1)^{2j}\,(2j{+}1)
      + \sum_{\beta} B_\beta(\chi)\, f(\beta, j)
    \Bigr],
\end{equation}
where the sum runs over all distinct half-angles
$\beta \in (0,\pi)$ of non-central conjugacy classes
in~$2H$, and
$B_\beta(\chi)
= \sum_{\substack{g \in 2H\\
\text{half-angle}(g)=\beta}}
\overline{\chi(g)}$
is the $\overline{\chi}$-weight at half-angle~$\beta$.
(For integer~$j$, one may group $\beta$ and
$\pi{-}\beta$ since
$f(\pi{-}\beta,j) = f(\beta,j)$; this fails at
half-integer~$j$ where
$f(\pi{-}\beta,j) = -f(\beta,j)$.)

\subsection{The divisibility condition}
\label{sec:divisibility}

Table~\ref{tab:spectral} records the scalar state count
for the principal configuration spaces.

\begin{table}[ht]
\centering
\begin{tabular}{@{}lcccc@{}}
  \toprule
  Space & $|2H|$ & $\kappa$ & $j_1$
        & $j^*$ \\
  \midrule
  $\SO(3)/T$   &   24 & 3 & $2$
    & $5$   \\
  $\SO(3)/O$   &   48 & 2 & $3$
    & $11$  \\
  $\SO(3)/I$   &  120 & 1 & $6$
    & $29$ \\
  \midrule
  $\SO(3)/D_3$ &   12 & 4 & $1$
    & $2$   \\
  \bottomrule
\end{tabular}
\caption{Scalar quantization data for the principal
  configuration spaces with $|H|$ even.
  $\kappa$: number of scalar sectors.
  $j_1$: lowest $j > 0$ with $m(\chi, j) \geq 1$
  for some scalar~$\chi$.
  $j^* = |H|/2 - 1$: spectral threshold~\eqref{eq:threshold}.
  Top three groups satisfy~(B); $D_3$ does not
  (\S\ref{sec:reentrance}).}
\label{tab:spectral}
\end{table}

For $H \subset \SO(3)$ finite with $|H|$ even, the
\emph{spectral threshold} is
\begin{equation}
  \label{eq:threshold}
  j^* = |H|/2 - 1.
\end{equation}
This is the unique value making $\dim V_j + \dim V_{j^*-j} = |H|$
for all~$j$.

\begin{definition}[Divisibility condition]
  \label{def:divisibility}
  The group $H \subset \SO(3)$ satisfies the
  \emph{divisibility condition}~(B) if
  $|H|/\ord(h)$ is even for every
  $h \in H \setminus \{e\}$.
\end{definition}

\noindent
Three equivalent formulations are useful:
\begin{enumerate}
  \item[\textup{(B$_1$)}] \emph{Element-wise:}
    $|H|/\ord(h)$ is even for every $h \neq e$.
  \item[\textup{(B$_2$)}] \emph{Exponent:}
    $\exp(H) \mid |H|/2$, where
    $\exp(H) = \lcm\{\ord(h) : h \in H\}$.
  \item[\textup{(B$_3$)}] \emph{$2$-adic:}
    $v_2(\ord(h)) < v_2(|H|)$ for every $h \neq e$.
\end{enumerate}
The equivalences are immediate:
(B$_1$)$\,\Leftrightarrow\,$(B$_2$) by taking the least common
multiple, and (B$_1$)$\,\Leftrightarrow\,$(B$_3$) by definition
of the $2$-adic valuation.
When $|H|$ is odd, $j^*$ is not an integer and~(B) fails
vacuously; such groups are excluded throughout.

% ============================================================
\section{Spectral reflection theorem}
\label{sec:duality}
% ============================================================

\subsection{Mode antisymmetry}
\label{sec:antisymmetry}

We use only that finite subgroups of $\SO(3)$ consist of
rotations of finite order, so each non-central conjugacy class
has half-angle $k'\pi/d$ with $d \mid |H|$.

\begin{lemma}[Mode antisymmetry]
  \label{lem:antisymmetry}
  If~\textup{(B)} holds, then for every non-central
  conjugacy class with half-angle~$\alpha$:
  \begin{equation}
    \label{eq:mode-antisymmetry}
    f(\alpha, j) + f(\alpha, j^*{-}j) = 0
    \qquad\text{for all } j.
  \end{equation}
\end{lemma}

\begin{proof}
  Set $K = |H|/d$ where $d$ is the $\SO(3)$ order
  of the class.
  An element of order~$d$
  is a rotation by angle $2\pi k'/d$ with
  $\gcd(k', d) = 1$, so its $\SU(2)$ lift has
  half-angle $\alpha = k'\pi/d$.
  Since $2j^*{+}1 = |H| - 1$:
  \[
    (2(j^*{-}j){+}1)\alpha
    = (|H|{-}1{-}2j)\,\frac{k'\pi}{d}
    = k'K\pi - (2j{+}1)\alpha.
  \]
  The identity
  $\sin(m\pi - x) = (-1)^{m+1}\sin(x)$ gives
  $f(\alpha, j^*{-}j) = (-1)^{k'K+1}\,f(\alpha, j)$.
  Condition~(B) requires $K$ even, so $k'K$ is even
  and $(-1)^{k'K+1} = -1$.
\end{proof}

\subsection{Exponent saturation}
\label{sec:exponent}

\begin{lemma}[Exponent saturation]
  \label{lem:exponent}
  If~\textup{(B)} holds, then $\exp(H) = |H|/2$.
\end{lemma}

\begin{proof}
  Under~(B), every $h \neq e$ has
  $\ord(h) \mid |H|/2$, so $|H| \ge 4$.
  Set $P = \exp(H) = \lcm\{\ord(h) : h \in H\}$.
  Each mode $f(\pi/d, j) = \sin((2j{+}1)\pi/d)/\sin(\pi/d)$
  has period~$d$ in~$j$, since
  $(2(j{+}d){+}1)\pi/d = (2j{+}1)\pi/d + 2\pi$.
  Because $P$ is divisible by every~$d$ that appears, the
  non-central contribution
  $N(j) = \sum_{g \neq \pm 1} \chi_j(g)$
  satisfies $N(j{+}P) = N(j)$.
  The central contribution to $m(A_0, j)$ is
  $C(j) = 2(2j{+}1)$, so $C(j{+}P) - C(j) = 4P$.
  The full difference is
  \[
    m(A_0,\, j{+}P) - m(A_0,\, j)
    = \frac{4P}{2|H|}
    = \frac{2P}{|H|}.
  \]
  Write $|H|/2 = kP$ with $k \in \ZZ_{\ge 1}$
  (possible since $P \mid |H|/2$ from~(B$_2$)).
  Then $m(A_0, j{+}P) - m(A_0, j) = 1/k$ for all
  integer~$j \ge 0$.
  Since $|\chi_j(g)| = |\sin((2j{+}1)\alpha_g)/\sin\alpha_g|
  \le 1/|\sin\alpha_g|$ for each $g \neq \pm 1$
  (with $\alpha_g$ the half-angle, fixed and nonzero),
  the non-central sum is bounded independently of~$j$.
  Hence $m(A_0, j) \ge (2(2j{+}1) - C_0)/|2H|$ for
  a constant~$C_0$, and $m(A_0, j) \to \infty$.
  The telescoping sum
  $m(A_0, nP) = m(A_0, 0) + n/k$ must be a
  non-negative integer for all~$n \ge 0$, forcing
  $1/k \in \ZZ_{>0}$.
  Since $k \ge 1$, the only possibility is $k = 1$,
  giving $P = |H|/2$.
\end{proof}

\noindent
Equivalently, (B) holds if and only if $\exp(H) = |H|/2$.
The threshold $j^* = P - 1 = |H|/2 - 1$ is thus uniquely
determined by~(B); no smaller threshold supports reflection.

\subsection{The spectral reflection theorem}
\label{sec:duality-theorem}

We now prove Theorem~\ref{thm:main} from the introduction.

\begin{theorem}[Spectral reflection]
  \label{thm:duality}
  For $H \subset \SO(3)$ finite with $|H|$ even, the
  following are equivalent:
  \begin{enumerate}
    \item[\textup{(B)}] The divisibility condition holds.
    \item[\textup{(C)}] For every scalar character~$\chi$
      and every $j$ in $[0, j^*]$ of the parity matching
      $\chi$
      \textup{(}integer if tensorial, half-integer if
      spinorial\textup{)}:
      \begin{equation}
        \label{eq:duality}
        m(\chi, j) + m(\chi, j^* - j) = 1.
      \end{equation}
  \end{enumerate}
\end{theorem}

\begin{proof}
  \textup{(B) $\Rightarrow$ (C).}
  By Lemma~\ref{lem:antisymmetry}, every non-central
  mode satisfies
  $f(\alpha, j) + f(\alpha, j^*{-}j) = 0$.
  Since the $\overline{\chi}$-weights
  in~\eqref{eq:mode-decomposition} are constants
  for each class, the entire non-central contribution
  to $m(\chi, j) + m(\chi, j^*{-}j)$ vanishes.
  The two central elements contribute
  \begin{align*}
    \frac{1}{|2H|}\Bigl[
      &\bigl((2j{+}1) + (2(j^*{-}j){+}1)\bigr) \\
      +\;\chi(-1)\bigl(
        &(-1)^{2j}(2j{+}1)
        + (-1)^{2(j^*-j)}(2(j^*{-}j){+}1)
      \bigr)
    \Bigr].
  \end{align*}
  The first line equals $|H|$ by definition of~$j^*$.
  For the second line:
  $(-1)^{2(j^*-j)} = (-1)^{|H|-2-2j}
  = (-1)^{|H|}(-1)^{2j}$,
  and since $|H|$ is even, $(-1)^{|H|} = 1$.
  The second line becomes
  $\chi(-1)\,(-1)^{2j}\,|H|$.
  The full central contribution is
  \begin{equation}
    \label{eq:central}
    \frac{1}{2}
    \bigl[1 + \chi(-1)\,(-1)^{2j}\bigr],
  \end{equation}
  which equals~$1$ when $\chi(-1) = (-1)^{2j}$
  (matching parity) and~$0$ otherwise.

  \medskip
  \textup{(C) $\Rightarrow$ (B).}
  Contrapositive.
  Suppose (B) fails: there exists $h \in H \setminus \{e\}$
  with $K = |H|/\mathrm{ord}(h)$ odd.
  For a non-central element $g \in 2H$ projecting to~$h$
  of $\SO(3)$ order~$d$,
  \begin{equation}
    \label{eq:chi-jstar}
    \chi_{j^*}(g)
    = \frac{\sin\bigl((|H|{-}1)\pi/d\bigr)}
           {\sin(\pi/d)}\,.
  \end{equation}
  Since $d \mid |H|$, write $(|H|{-}1)/d = K - 1/d$.
  Then $(|H|{-}1)\pi/d = K\pi - \pi/d$, so
  $\sin((|H|{-}1)\pi/d) = (-1)^{K+1}\sin(\pi/d)$,
  giving $\chi_{j^*}(g) = (-1)^{K+1} = +1$ when $K$ is odd.

  Now compute $m(A_0, j^*)$ from~\eqref{eq:multiplicity}.
  Each non-identity $h \in H$ lifts to a pair
  $g_h, -g_h$ in~$2H$.
  Since $j^*$ is an integer,
  $\chi_{j^*}(-g_h) = \chi_{j^*}(g_h)$, and the pair
  contributes $2\chi_{j^*}(g_h)$: this is $+2$ when
  $|H|/\mathrm{ord}(h)$ is odd, and $-2$ when even.
  Let $S_{\mathrm{odd}} = \{h \in H\setminus\{e\} : |H|/\mathrm{ord}(h)
  \text{ is odd}\}$.
  The central elements contribute
  $2(2j^*{+}1) = 2(|H|{-}1)$.
  Writing $S_{\mathrm{even}} = (H\setminus\{e\})
  \setminus S_{\mathrm{odd}}$:
  \[
    |2H|\,m(A_0, j^*)
    = 2(|H|{-}1)
    + 2|S_{\mathrm{odd}}|
    - 2|S_{\mathrm{even}}|.
  \]
  Since $|S_{\mathrm{odd}}| + |S_{\mathrm{even}}| = |H|{-}1$,
  this gives
  \begin{equation}
    \label{eq:m-jstar}
    m(A_0, j^*)
    = \frac{2}{|H|}
      |S_{\mathrm{odd}}|.
  \end{equation}
  Since $m(A_0,j^*) \in \ZZ_{\ge 0}$,
  $|S_{\mathrm{odd}}|$ is a non-negative multiple
  of~$|H|/2$.
  Since $|S_{\mathrm{odd}}| \le |H|{-}1$, the only
  possibilities are $0$ and~$|H|/2$.
  Since (B) fails, $S_{\mathrm{odd}} \neq \emptyset$,
  so $|S_{\mathrm{odd}}| = |H|/2$ and
  $m(A_0, j^*) = 1$.
  Combined with $m(A_0, 0) = 1$:
  $m(A_0, 0) + m(A_0, j^*) = 2 \neq 1$.
\end{proof}

\subsection{Consequences}
\label{sec:consequences}

\begin{corollary}[Binary occupation]
  \label{cor:binary}
  Under~\textup{(B)}, $m(\chi, j) \in \{0,1\}$ for every
  scalar~$\chi$ and every $j \le j^*$ of matching parity.
\end{corollary}

\begin{proof}
  From $m(\chi,j) + m(\chi,j^*{-}j) = 1$ with both terms
  non-negative integers, each is $0$ or~$1$.
\end{proof}

\begin{corollary}[Regular-representation tiling]
  \label{cor:reg-rep}
  Under~\textup{(B)}, for every integer $j \in [0,\, j^*]$:
  \begin{equation}
    \label{eq:reg-rep-tiling}
    \mathrm{Res}\,(V_j \oplus V_{j^*-j})\big|_{2H}
    \cong \rho_{\mathrm{reg}}(H),
  \end{equation}
  where $\rho_{\mathrm{reg}}(H) = \bigoplus_{\rho} (\dim\rho)\,\rho$
  is the regular representation of~$H$.
\end{corollary}

\begin{proof}
  For any irreducible representation~$\rho$ of~$2H$
  of matching parity ($\rho(-1) = I$ at integer~$j$,
  $\rho(-1) = -I$ at half-integer~$j$), the
  multiplicity $m(\rho,j) = \frac{1}{|2H|}
  \sum_g \overline{\chi_\rho(g)}\,\chi_j(g)$
  satisfies the same paired cancellation:
  Lemma~\ref{lem:antisymmetry} kills the non-central
  terms (it depends only on the half-angle, not on
  the character weighting), and the central contribution
  gives
  $\frac{1}{|2H|}[\dim\rho \cdot |H|
  + \chi_\rho(-1)\,(-1)^{2j}\,\dim\rho \cdot |H|]
  = \dim\rho$
  when the parity matches.
  Hence
  \[
    m(\rho, j) + m(\rho, j^*{-}j) = \dim\rho.
  \]
  At integer~$j$, only representations with
  $\rho(-1) = I$ contribute; these are precisely the
  irreps of~$2H$ that factor through $H = 2H/\{\pm 1\}$.
  Summing $(\dim\rho)\,\rho$ over matching-parity irreps
  gives $\rho_{\mathrm{reg}}(H)$.
\end{proof}

\begin{corollary}[Periodic staircase]
  \label{cor:staircase}
  Under~\textup{(B)}, for every scalar~$\chi$ and
  every $j \ge 0$ of matching parity:
  \begin{equation}
    \label{eq:staircase}
    m(\chi,\, j + P) = m(\chi,\, j) + 1,
  \end{equation}
  where $P = |H|/2 = j^* + 1$.
\end{corollary}

\begin{proof}
  By Lemma~\ref{lem:exponent}, every non-central mode
  $f(\alpha, j)$ has period dividing $P$ in~$j$.
  Hence the non-central sum in~\eqref{eq:mode-decomposition}
  satisfies $N(j{+}P) = N(j)$.
  The central contribution grows by
  $2(2P)/|2H| = 2P/|H| = 1$ per period,
  giving~\eqref{eq:staircase}.
  The identity holds for all $j \ge 0$, including above threshold.
\end{proof}

\begin{corollary}[Threshold sharpness]
  \label{cor:threshold}
  Under~\textup{(B)}, for every scalar~$\chi$ of
  matching parity:
  \begin{equation}
    \label{eq:threshold-sharp}
    m(\chi,\, j^*{+}1) = m(\chi, 0) + 1.
  \end{equation}
  In particular, $m(A_0,\, j^*{+}1) = 2$.
\end{corollary}

\begin{proof}
  Set $j = 0$ in~\eqref{eq:staircase}:
  $m(\chi, P) = m(\chi, 0) + 1$.
  Since $P = j^*{+}1$, the result follows.
  For $A_0$: $m(A_0, 0) = 1$, so
  $m(A_0, j^*{+}1) = 2$.
  For characters with $m(\chi, 0) = 0$
  (e.g., $A_1$ in~$O$),
  $m(\chi, j^*{+}1) = 1$: the level appears
  for the first time above threshold.
\end{proof}

\begin{remark}[Periodic decomposition]
  \label{rem:periodic}
  Corollaries~\ref{cor:staircase}
  and~\ref{cor:binary} together yield a
  decomposition
  $m(\chi, j) = \lfloor j/P \rfloor + \varepsilon(\chi, j)$,
  where $\varepsilon(\chi, \cdot)$ is periodic with
  period~$P$, takes values in~$\{0,1\}$, and satisfies
  $\varepsilon(\chi, j) + \varepsilon(\chi, P{-}1{-}j) = 1$
  (the reflection identity restricted to the first period).
  This reformulates the spectral structure as a linear
  drift with a periodic binary residual.
\end{remark}

% ============================================================
\section{Topological obstruction and symmetry breaking}
\label{sec:obstruction}
% ============================================================

\subsection{The octahedral rotor}
\label{sec:octahedral}

For $H = O$ (order~$24$), the binary cover
$2O$ has order~$48$, abelianization~$\ZZ_2$,
and two scalar sectors: the trivial character~$A_0$
and the sign character~$A_1$ (both tensorial).
The threshold is $j^* = 11$.

Table~\ref{tab:octahedral} displays the scalar
multiplicities at all sub-threshold levels.
The reflection $m(\chi, j) + m(\chi, 11{-}j) = 1$
holds in every row.

\begin{table}[ht]
\centering
\begin{tabular}{@{}ccccc@{}}
  \toprule
  $j$ & $m(A_0, j)$ & $m(A_1, j)$
    & $m(A_0, 11{-}j)$ & $m(A_1, 11{-}j)$ \\
  \midrule
  $0$  & $1$ & $0$ & $0$ & $1$ \\
  $1$  & $0$ & $0$ & $1$ & $1$ \\
  $2$  & $0$ & $0$ & $1$ & $1$ \\
  $3$  & $0$ & $1$ & $1$ & $0$ \\
  $4$  & $1$ & $0$ & $0$ & $1$ \\
  $5$  & $0$ & $0$ & $1$ & $1$ \\
  $6$  & $1$ & $1$ & $0$ & $0$ \\
  $7$  & $0$ & $1$ & $1$ & $0$ \\
  $8$  & $1$ & $0$ & $0$ & $1$ \\
  $9$  & $1$ & $1$ & $0$ & $0$ \\
  $10$ & $1$ & $1$ & $0$ & $0$ \\
  $11$ & $0$ & $1$ & $1$ & $0$ \\
  \bottomrule
\end{tabular}
\caption{Scalar multiplicities for $H = O$,
  $j^* = 11$.
  Both characters are tensorial (integer~$j$ only).
  The reflection $m(\chi, j) + m(\chi, 11{-}j) = 1$
  holds in every row.}
\label{tab:octahedral}
\end{table}

\noindent
For $A_0$, the
mode decomposition~\eqref{eq:mode-decomposition}
has three non-central modes:
\begin{equation}
  \label{eq:collapsed-O}
  S_O(j)
  = 12\,f(\tfrac{\pi}{4}, j)
  + 16\,f(\tfrac{\pi}{3}, j)
  + 18\,f(\tfrac{\pi}{2}, j),
\end{equation}
corresponding to $\SO(3)$~orders
$d = 4, 3, 2$ with $K = |O|/d = 6, 8, 12$---all even,
confirming~(B).

\subsection{The topological obstruction}
\label{sec:topological}

Both octahedral sectors are tensorial---this is forced
by a topological obstruction.
Under a $2\pi$ rotation, a wave function in sector~$\chi$
acquires the phase $\chi(-1)$.
If $-1$ lies in the commutator subgroup
$[G,G]$ of $G = 2H$, then $\chi(-1) = 1$ for every
one-dimensional character: no spinorial sector exists.
\begin{equation}
  \label{eq:obstruction-A}
  \textup{(A)}\qquad -1 \in [G,G].
\end{equation}

\begin{proposition}[Topological link]
  \label{prop:A-iff-B}
  For $G = 2H \subset \SU(2)$, condition~\textup{(A)}
  is equivalent to the divisibility
  condition~\textup{(B)}.
\end{proposition}

\begin{proof}
  \emph{(A)$\,\Rightarrow\,$(B).}\;
  If $[a,b] = -1$ for some $a,b\in 2H$,
  then $ab = -ba$.
  Since $a^2$ commutes with~$b$ and lies in
  $\{\pm 1\}$, with $a^2 = +1$ giving $[a,b] = 1$,
  we have $a^2 = b^2 = -1$.
  Hence $a,b$ are pure imaginary unit quaternions and
  $\{\pm 1,\pm a,\pm b,\pm ab\}\cong Q_8\subset 2H$.
  Projecting gives $V_4\subset H$, so the
  Sylow 2-subgroup~$S$ of~$H$ is non-cyclic.
  Since $S$ is a finite 2-subgroup of $\SO(3)$, the
  Burnside orbit-counting equation on~$S^2$
  forces $S \cong D_{2^{a-1}}$
  (dihedral of order~$2^a$), with exponent~$2^{a-1}$.
  For any $h \in H$, writing $\ord(h) = 2^b m'$ ($m'$ odd),
  the element $h^{m'}$ has 2-power order~$2^b$ and lies
  in a Sylow 2-subgroup, so $b \le a{-}1$.
  Thus $\nu_2(\ord h) < \nu_2(|H|)$ for all
  $h\neq e$, which is~(B${}_3$).

  \emph{(B)$\,\Rightarrow\,$(A).}\;
  Write $|H|=2^a m$ ($m$ odd, $a\ge 1$).
  \emph{Step~1:} If the Sylow 2-subgroup~$S$ were cyclic,
  it would contain an element of order~$2^a$,
  violating~(B${}_3$).  So $S$ is non-cyclic,
  with $|S|\ge 4$ since $\exp(H) = |H|/2$ forces $|H| \ge 4$.
  \emph{Step~2:} The generalised quaternion group
  $Q_{2^n}$ ($n\ge 3$) admits no faithful homomorphism
  to~$\SO(3)$: its presentation forces
  $\varphi(a)$ and $\varphi(b)$ to share a rotation axis
  (both $\varphi(b)^2$ and $\varphi(a^{2^{n-2}})$
  are the same $\pi$-rotation), but $bab^{-1}=a^{-1}$
  with $\ord(a)\ge 4$ requires non-commutativity,
  a contradiction.
  \emph{Step~3:} By Burnside's theorem~\cite{Gorenstein1980},
  a finite 2-group of order~$\ge 4$ that is
  neither cyclic nor generalised quaternion contains
  $V_4 = \ZZ_2\times\ZZ_2$.
  \emph{Step~4:} The three non-identity elements of
  $V_4\subset H$ are $\pi$-rotations with mutually
  perpendicular axes; their $\SU(2)$ lifts are pure
  imaginary and pairwise anticommute, generating
  $Q_8\subset 2H$.
  Then $[a,b]=-1$, giving~(A).
\end{proof}

\begin{table}[ht]
\centering
\begin{tabular}{@{}lcccc@{}}
  \toprule
  $H \subset \SO(3)$ & (B)? & $-1 \in [2H,2H]$?
    & $2\pi$ holonomy & spinorial? \\
  \midrule
  $C_n$ & No & No & $-1$ allowed & Yes \\
  $D_n$ ($n$ odd) & No & No & $-1$ allowed & Yes \\
  $D_n$ ($n$ even) & Yes & Yes & $+1$ forced & No \\
  $T,O,I$ & Yes & Yes & $+1$ forced & No \\
  \bottomrule
\end{tabular}
\caption{Obstruction and divisibility for all finite
  $H \subset \SO(3)$.
  When (B) holds, every scalar sector has $\chi(-1) = +1$.
  Even dihedral groups satisfy~(B) with
  abelianization $\ZZ_2 \times \ZZ_2$
  (four tensorial scalar sectors).}
\label{tab:obstruction-summary}
\end{table}

\subsection{Symmetry breaking: re-entrance and classification}
\label{sec:reentrance}

When the rotation group is lowered from $H$ to a subgroup
$K \subset H$, two spectral effects occur:
(1)~\emph{densification}: levels forbidden by $H$-symmetry
become allowed under~$K$, and the spectral gap decreases;
(2)~\emph{re-entrance}: spinorial sectors may reappear.
The mechanism is arithmetic: re-entrance is controlled by
whether the subgroup chain eliminates~$Q_8$.

\medskip
\begin{center}
\begin{tabular}{@{}lcc@{}}
  \toprule
  & Trigonal ($O \to D_3$) & Tetragonal ($O \to D_4$) \\
  \midrule
  $Q_8 \subset 2K$?     & No         & Yes        \\
  1D characters           & $2 \to 4$  & $2 \to 4$  \\
  \quad of which spinorial & $0 \to 2$ (new) & $0 \to 0$ \\
  Smallest tensorial gap  & $2B$ ($\chi_2$, $j{=}1$)
                          & $6B$ ($A_0$, $j{=}2$) \\
  Spinorial ground state  & $j{=}3/2$, $\Delta{=}3.75B$ & --- \\
  \bottomrule
\end{tabular}
\end{center}

\medskip
\noindent
We develop the trigonal case.
The group $2D_3 = \mathrm{Dic}_3 \subset \SU(2)$ has
order~$12$, with abelianization~$\ZZ_4$.
Its four one-dimensional characters are:
$\chi_0$~(T): trivial; $\chi_2$~(T): sign of~$D_3$
($\chi_2(-1) = +1$); $\chi_1$~(S), $\chi_3$~(S): spinorial
($\chi(-1) = -1$).
The character table, with classes labelled by
half-angle~$\alpha$:

\begin{center}
\begin{tabular}{@{}lcccccc@{}}
  \toprule
  class
    & $\{1\}$ & $\{a,a^5\}$ & $\{a^2,a^4\}$
    & $\{-1\}$ & $\{b,a^2b,a^4b\}$ & $\{ab,a^3b,a^5b\}$ \\
  size & $1$ & $2$ & $2$ & $1$ & $3$ & $3$ \\
  $\alpha$ & $0$ & $\pi/3$ & $2\pi/3$ & $\pi$
    & $\pi/2$ & $\pi/2$ \\
  \midrule
  $\chi_0$ & $1$ & $1$  & $1$  & $1$  & $1$  & $1$  \\
  $\chi_1$ & $1$ & $-1$ & $1$  & $-1$ & $i$  & $-i$ \\
  $\chi_2$ & $1$ & $1$  & $1$  & $1$  & $-1$ & $-1$ \\
  $\chi_3$ & $1$ & $-1$ & $1$  & $-1$ & $-i$ & $i$  \\
  \bottomrule
\end{tabular}
\end{center}

\noindent
The multiplicity table:

\begin{center}
\begin{tabular}{@{}ccccc@{}}
  \toprule
  $j$ & $\chi_0$ (T) & $\chi_1$ (S) & $\chi_2$ (T) & $\chi_3$ (S) \\
  \midrule
  $0$   & $1$ & ---  & $0$ & ---  \\
  $1/2$ & --- & $0$  & --- & $0$  \\
  $1$   & $0$ & ---  & $1$ & ---  \\
  $3/2$ & --- & $1$  & --- & $1$  \\
  $2$   & $1$ & ---  & $0$ & ---  \\
  $5/2$ & --- & $1$  & --- & $1$  \\
  \bottomrule
\end{tabular}
\end{center}

\noindent
The tensorial reflection fails:
$m(\chi_0, 0) + m(\chi_0, 2) = 1 + 1 = 2 \neq 1$.

\begin{table}[ht]
\centering
\begin{tabular}{@{}cccccc@{}}
  \toprule
  Sector & Type & Symmetry & Ancestor & $j_1$ & $\Delta/B$ \\
  \midrule
  $A_0$ & T & $O$   & ---       & $4$   & $20$ \\
  $A_1$ & T & $O$   & ---       & $3$   & $12$ \\
  \midrule
  $\chi_0$ & T & $D_3$ & $A_0$  & $2$   & $6$ \\
  $\chi_2$ & T & $D_3$ & $A_1$  & $1$   & $2$ \\
  $\chi_1$ & S & $D_3$ & none   & $3/2$ & $3.75$ \\
  $\chi_3$ & S & $D_3$ & none   & $3/2$ & $3.75$ \\
  \bottomrule
\end{tabular}
\caption{Spectral structure before and after $O \to D_3$.
  (T)~tensorial, (S)~spinorial.
  Top: $\SO(3)/O$ (obstruction active).
  Bottom: $\SO(3)/D_3$ (obstruction lifts, spinorial sectors
  appear with no ancestor in~$O$).}
\label{tab:densification}
\end{table}

The trigonal and tetragonal examples are instances of a general
principle.

\begin{proposition}[Lifting criterion]
  \label{prop:lifting}
  Let $H \subset \SO(3)$ satisfy~\textup{(B)}, and let
  $K \subset H$ be a subgroup.
  Spinorial scalar sectors exist on $\SO(3)/K$ if and only
  if $Q_8 \not\subset 2K$.
\end{proposition}

\begin{proof}
  A spinorial character of~$2K$ exists iff
  $-1 \notin [2K, 2K]$, which by
  Proposition~\ref{prop:A-iff-B} (applied to~$K$)
  is equivalent to $Q_8 \not\subset 2K$.
\end{proof}

\begin{table}[ht]
\centering
\begin{tabular}{@{}lcccc@{}}
  \toprule
  Channel & $|2K|$ & $Q_8 \subset 2K$?
    & spinorial? & mechanism \\
  \midrule
  $O$ (parent) & $48$ & Yes & No & $Q_8 \subset 2O$ \\
  $O \to T$    & $24$ & Yes & No & $Q_8 \subset 2T$ \\
  $O \to D_4$  & $16$ & Yes & No & $Q_8 \subset 2D_4$ \\
  $O \to D_2$  & $8$  & Yes & No & $Q_8 = 2D_2$ \\
  \midrule
  $O \to D_3$  & $12$ & No  & Yes & Lagrange: $8 \nmid 12$ \\
  $O \to C_4$  & $8$  & No  & Yes & $\ZZ_8$ abelian \\
  $O \to C_3$  & $6$  & No  & Yes & Lagrange: $8 \nmid 6$ \\
  $O \to C_2$  & $4$  & No  & Yes & $\ZZ_4$ abelian \\
  \bottomrule
\end{tabular}
\caption{Symmetry-breaking phase diagram for the octahedral
  rotor.
  Above the midrule: obstruction persists ($Q_8 \subset 2K$).
  Below: obstruction lifts, spinorial sectors reappear.}
\label{tab:lifting}
\end{table}

\noindent
The criterion is monotone: once $Q_8$ is lost in a subgroup
chain $H \supset K \supset K'$, it cannot reappear.
This completes the classification: for every finite
$H \subset \SO(3)$ and every subgroup $K \subset H$,
the spectral structure on $\SO(3)/K$ is determined
by Theorem~\ref{thm:duality} when~(B) holds for~$K$,
and by the explicit failure mode when~(B) fails.

\subsection{Spinorial gap}
\label{sec:spinorial-gap}

Proposition~\ref{prop:lifting} characterises \emph{when}
spinorial sectors exist.
We now determine \emph{where} they first appear.

\begin{proposition}[Spinorial gap]
  \label{prop:spinorial-gap}
  Let $K \subset \SO(3)$ with $Q_8 \not\subset 2K$.
  \begin{enumerate}
  \item If $K \cong C_n$, the lowest spinorial level is
    $j = \tfrac{1}{2}$, with gap
    $\Delta_{\mathrm{spin}} = \tfrac{3}{4}\,B$.
  \item If $K \cong D_n$ with $n$ odd, the lowest spinorial
    level is $j_{\min}^{\mathrm{spin}} = n/2$, with gap
    \[
      \Delta_{\mathrm{spin}}
      = B\,\frac{n}{2}\Bigl(\frac{n}{2}+1\Bigr)
      = \frac{n(n+2)}{4}\,B.
    \]
  \end{enumerate}
  The spinorial gap grows quadratically
  in~$n$ for dihedral subgroups.
\end{proposition}

\begin{proof}
  For $C_n$: $V_{1/2}|_{\ZZ_{2n}}$ contains the character
  $\chi_1$ with $\chi_1(-1) = -1$, so $m(\chi_1, \tfrac{1}{2}) \geq 1$.

  For $D_n$ ($n$ odd): $2D_n \cong \mathrm{Dic}_n$.
  Every spinorial character restricts to $\chi_n$ on the
  cyclic subgroup $\ZZ_{2n} = \langle a \rangle$.
  The weights of $V_j|_{\ZZ_{2n}}$ are
  $\{2m \bmod 2n : m = -j, \ldots, j\}$.
  The character~$\chi_n$ appears iff $2m \equiv n \pmod{2n}$,
  i.e.\ $m = n/2$, which requires $j \geq n/2$.
  At $j = n/2$ the weight is realised.
\end{proof}

\begin{remark}
  The $D_3$ case ($j_{\min}^{\mathrm{spin}} = 3/2$,
  $\Delta = 3.75B$) is the lowest dihedral instance.
  For $D_5$: $\Delta = 8.75B$; for $D_7$: $\Delta = 15.75B$.
  Existence of spinorial sectors does not imply
  low-energy spinorial excitations.
\end{remark}

\subsection{Sector-selective duality}
\label{sec:sector-selective}

Although Theorem~\ref{thm:duality} fails globally for $D_n$
with $n$ odd, the reflection persists in the spinorial sector.
The failure of~(B) is carried entirely by half-turn classes
($d = 2$), which are spectrally invisible at half-integer~$j$.

\begin{lemma}[Spectral invisibility]
  \label{lem:invisibility}
  A conjugacy class of $\SO(3)$ order~$d$ makes zero contribution
  to the multiplicity formula~\eqref{eq:multiplicity}
  at every $j$ with $d \mid 2j{+}1$.
  In particular, half-turn classes ($d = 2$) are invisible
  at all half-integer~$j$.
\end{lemma}

\begin{proof}
  The mode function $f(\pi/d, j) = \sin((2j{+}1)\pi/d)/\sin(\pi/d)$
  vanishes whenever $d \mid 2j{+}1$.
  For $d = 2$: at half-integer~$j$, $2j{+}1$ is even,
  so $\sin((2j{+}1)\pi/2) = 0$.
\end{proof}

\begin{proposition}[Sector-selective duality]
  \label{prop:sector-selective}
  Let $H = D_n$ with $n$ odd.
  Then condition~\textup{(B)} fails, but the spectral
  reflection holds for all spinorial characters:
  $m(\chi,j) + m(\chi,j^*{-}j) = 1$ for every spinorial
  $\chi$ and every half-integer $j \in [\tfrac{1}{2},\,
  j^*{-}\tfrac{1}{2}]$.
  The reflection fails for tensorial characters.
\end{proposition}

\begin{proof}
  The only classes violating~(B) in $D_n$ ($n$ odd) are
  the half-turns ($d = 2$, $|H|/d = n$ odd).
  By Lemma~\ref{lem:invisibility}, these are invisible at
  half-integer~$j$.
  The remaining non-central classes satisfy $|H|/d$ even, so
  Lemma~\ref{lem:antisymmetry} applies class-by-class
  and their contributions cancel in the paired sum.
  The central contribution equals~$1$ for spinorial
  $\chi$ at half-integer~$j$, since
  $\chi(-1)\,(-1)^{2j} = (-1)(-1) = 1$.

  For tensorial characters ($\chi(-1) = 1$, integer~$j$),
  the half-turn classes are \emph{not} invisible:
  $f(\pi/2, j) = (-1)^j \neq 0$.
  Since $K = n$ is odd,
  $f(\pi/2,\, j^*{-}j) = +f(\pi/2, j)$
  (the contribution doubles rather than cancels),
  and the reflection sum becomes~$0$ or~$2$.
\end{proof}

\begin{remark}[Generalized selective duality]
  \label{rem:generalized-selective}
The proof uses only mode antisymmetry
(Lemma~\ref{lem:antisymmetry}) for classes satisfying~(B),
and spectral invisibility (Lemma~\ref{lem:invisibility})
for the remaining classes.
A scalar character~$\chi$ satisfies the reflection
whenever $\chi$ is immune to every (B)-violating class---that is,
whenever $f(\alpha_\beta, j) = 0$ for all~$j$ in the
support of~$\chi$ and every class~$\beta$ with
$|H|/\ord(\beta)$ odd.
Among finite subgroups of $\SO(3)$ with $|H|$ even
and~(B) failing, $D_n$ with $n$ odd is the only family
whose violating classes consist entirely of half-turns;
for $C_n$ with $n$ even, the violating elements include
rotations of order $n > 2$, which are not spectrally
invisible, and no selective duality occurs.
\end{remark}

% ============================================================
\section{Discussion}
\label{sec:discussion}
% ============================================================

Theorem~\ref{thm:duality} and Proposition~\ref{prop:A-iff-B}
together determine the sub-threshold spectral structure
for all finite $H \subset \SO(3)$ with $|H|$ even.
The partition is threefold:
\begin{enumerate}
\item[(i)] When~(B) holds ($D_n$ with $n$ even, $T$, $O$, $I$):
  the reflection identity, binary occupation, periodic
  staircase, and regular-representation tiling all hold;
  every scalar sector is tensorial.
\item[(ii)] When~(B) fails with $|H|$ even ($D_n$ with
  $n$ odd): the reflection fails for tensorial characters
  but persists in spinorial sectors by selective duality
  (Proposition~\ref{prop:sector-selective}).
\item[(iii)] When $|H|$ is odd ($C_n$ with $n$ odd):
  $j^*$ is not an integer and no reflection structure exists.
\end{enumerate}

\emph{Invariant-theory interpretation.}
The multiplicities $m(A_0, j)$ are the
coefficients of the Molien series
$M(t) = |2H|^{-1} \sum_{g \in 2H} \det(I - tg)^{-1}$
of the invariant ring $\CC[x,y]^{2H}$, evaluated at
degree~$2j$ (since $V_j = \Sym^{2j}(\CC^2)$).
For $2H \subset \SL(2,\CC)$,
$\CC[x,y]^{2H}$ is Gorenstein~\cite{Watanabe1974},
so $M(t)$ takes the rational form
$(1 + t^s)/\bigl((1-t^{d_1})(1-t^{d_2})\bigr)$
with palindromic numerator,
where $d_1, d_2$ are the degrees of a homogeneous system
of parameters and $s = d_1 + d_2 - 2$:
explicitly, $(d_1,d_2,s) = (6,8,12)$ for~$T$,
$(8,12,18)$ for~$O$, and $(12,20,30)$ for~$I$.
The palindromic rational form implies a functional symmetry
of~$M(t)$; combined with the fact that $m(A_0, j) \in \{0,1\}$
below degree~$|H| - 2$ (which itself follows from the
mode-antisymmetry argument), reflection for~$A_0$ follows.
The full identity for all scalar~$\chi$---and the equivalence
(B)$\,\Leftrightarrow\,$(C)---requires
Theorem~\ref{thm:duality}.

\emph{Physical context.}
For a rigid molecular rotor, nuclear spins transform
under~$H$, not its double cover~$2H$, so spinorial
sectors carry zero nuclear-spin statistical weight
$g_\chi = 0$~\cite{BunkerJensen2006}.
These sectors are quantum-mechanically allowed but
thermally dark in standard molecular spectroscopy.
The tensorial sectors, however, directly carry the
reflection signature: the binary occupation pattern
$m(\chi,j) \in \{0,1\}$ and the complementarity
$m(\chi,j) + m(\chi,j^*{-}j) = 1$ constrain which
rotational lines appear in each sector.
In engineered systems without nuclear-spin degrees of
freedom---such as atoms in optical lattices with
point-group potentials---all sectors are populated
and the full spectral structure is observable.

\emph{Scope and limitations.}
The mechanism relies on the linear dimension formula
$\dim V_j = 2j{+}1$; for groups of higher rank,
$\dim V_\lambda$ grows polynomially and the arithmetic
underpinning the reflection breaks down.
For $\SU(3)$, $\dim V_{(p,q)} = \tfrac{1}{2}(p{+}1)(q{+}1)(p{+}q{+}2)$
is quadratic; no affine involution on the weight lattice
produces a constant paired sum, and finite subgroups such as
$\Delta(27)$ exhibit no complementary pairing
(verified computationally).
A natural question is whether analogous results hold for
spherical space forms $S^3/\Gamma$~\cite{Wolf2011}.
The Laplacian on~$S^3$ decomposes under
$\SO(4) \cong (\SU(2) \times \SU(2))/\ZZ_2$, and
restriction to $\Gamma \subset \SU(2)$ acting on one factor
involves two-sided representation theory~\cite{Ikeda1980}:
eigenspaces are $V_\ell \otimes V_\ell$, so multiplicities
involve $\dim\Hom_\Gamma(V_\ell, V_\ell)$ rather than the
one-sided branching $\dim V_j^H$ used here.

\emph{Structural content.}
The arguments derive spectral structure directly from
group order and element orders, without enumerating
the polyhedral types.
The triple equivalence reduces to:
(B) holds $\Leftrightarrow$ the Sylow 2-subgroup
of~$H$ is non-cyclic $\Leftrightarrow$
$Q_8 \subset 2H$ $\Leftrightarrow$
$-1 \in [2H, 2H]$.

% ============================================================
\begin{thebibliography}{99}

\bibitem{Schulman1968}
  L.~S.~Schulman,
  ``A path integral for spin,''
  Phys.\ Rev.\ \textbf{176}, 1558 (1968).

\bibitem{Dowker1972}
  J.~S.~Dowker,
  ``Quantum mechanics and field theory on multiply connected
  and on homogeneous spaces,''
  J.\ Phys.\ A: Gen.\ Phys.\ \textbf{5}, 936 (1972).

\bibitem{Isham1984}
  C.~J.~Isham,
  ``Topological and global aspects of quantum theory,''
  in \emph{Relativity, Groups and Topology~II}
  (Les~Houches, 1983),
  ed.\ B.~S.~DeWitt and R.~Stora,
  North-Holland (1984), pp.~1059--1290.

\bibitem{LaidlawDeWitt1971}
  M.~G.~G.~Laidlaw and C.~M.~DeWitt,
  ``Feynman functional integrals for systems of
  indistinguishable particles,''
  Phys.\ Rev.\ D \textbf{3}, 1375 (1971).

\bibitem{ImboSudarshan1988}
  T.~D.~Imbo and E.~C.~G.~Sudarshan,
  ``Inequivalent quantizations and fundamentally
  perfect spaces,''
  Phys.\ Rev.\ Lett.\ \textbf{60}, 481 (1988).

\bibitem{Mermin1979}
  N.~D.~Mermin,
  ``The topological theory of defects in ordered media,''
  Rev.\ Mod.\ Phys.\ \textbf{51}, 591 (1979).

\bibitem{Giulini1995}
  D.~Giulini,
  ``On the selection rules for the configuration-space
  topology in quantum mechanics,''
  Helv.\ Phys.\ Acta \textbf{68}, 439 (1995).

\bibitem{Mouchet2021}
  A.~Mouchet,
  ``Path integrals in a multiply-connected configuration
  space (50 years after),''
  Found.\ Phys.\ \textbf{51}, 97 (2021).

\bibitem{Gorenstein1980}
  D.~Gorenstein,
  \emph{Finite Groups},
  2nd ed., Chelsea (1980).

\bibitem{BrockerTomDieck1985}
  T.~Br\"ocker and T.~tom~Dieck,
  \emph{Representations of Compact Lie Groups},
  Graduate Texts in Mathematics \textbf{98},
  Springer (1985).

\bibitem{BunkerJensen2006}
  P.~R.~Bunker and P.~Jensen,
  \emph{Molecular Symmetry and Spectroscopy},
  2nd ed., NRC Research Press (2006).

\bibitem{FinkelsteinRubinstein1968}
  D.~Finkelstein and J.~Rubinstein,
  ``Connection between spin, statistics, and kinks,''
  J.\ Math.\ Phys.\ \textbf{9}, 1762 (1968).

\bibitem{Hecht1960}
  K.~T.~Hecht,
  ``The vibration-rotation energies of tetrahedral
  XY$_4$ molecules: Part~I.\ Theory of spherical
  top molecules,''
  J.\ Mol.\ Spectrosc.\ \textbf{5}, 355 (1960).

\bibitem{Wolf2011}
  J.~A.~Wolf,
  \emph{Spaces of Constant Curvature},
  6th ed., AMS Chelsea (2011).

\bibitem{BalachandranEtAl1983}
  A.~P.~Balachandran, G.~Marmo, B.~S.~Skagerstam,
  and A.~Stern,
  ``Gauge symmetries and fibre bundles:
  Applications to particle dynamics,''
  Lecture Notes in Physics \textbf{188},
  Springer (1983).

\bibitem{Ikeda1980}
  A.~Ikeda,
  ``On the spectrum of a Riemannian manifold of positive
  constant curvature. II,''
  Osaka J.\ Math.\ \textbf{17}, 75 (1980).

\bibitem{Watanabe1974}
  K.~Watanabe,
  ``Certain invariant subrings are Gorenstein.\ I,''
  Osaka J.\ Math.\ \textbf{11}, 1--8 (1974).

\bibitem{AltmannHerzig1994}
  S.~L.~Altmann and P.~Herzig,
  \emph{Point-Group Theory Tables},
  2nd ed., Oxford University Press (1994).

\bibitem{Zhilinskii2001}
  B.~Zhilinski\'{\i},
  ``Symmetry, invariants, and topology in molecular
  models,''
  Phys.\ Rep.\ \textbf{341}, 85--171 (2001).

\end{thebibliography}

\end{document}
